\chapter{PENDAHULUAN}

\petunjuk{Gunakan bab ini untuk menjelaskan latar belakang, rumusan masalah, batasan masalah, tujuan, manfaat, metodologi singkat, dan sistematika penulisan. Hapus seluruh kotak petunjuk sebelum penyerahan akhir.}

\section{Latar Belakang}

Tuliskan latar belakang secara runtut dari konteks umum menuju masalah khusus yang akan diselesaikan. Setiap paragraf menggunakan Times New Roman 12 pt, spasi 1,5, rata kanan-kiri, dan indentasi baris pertama 1 cm. Hindari paragraf yang hanya terdiri dari satu kalimat.

\subsection{Contoh Anak Subbab}

Anak subbab menggunakan penomoran tiga tingkat, misalnya 1.1.1. Anak subbab dicetak tebal dan dimasukkan ke dalam Daftar Isi.

\subsubsection*{Contoh Cucu Subbab}

Cucu subbab dicetak tebal tanpa nomor. Cucu subbab tidak dimasukkan ke dalam Daftar Isi. Gunakan tingkat ini hanya bila uraian benar-benar memerlukan pembagian tambahan di bawah anak subbab.

\section{Rumusan Masalah}

Rumusan masalah harus konsisten dengan tujuan dan kesimpulan. Apabila perlu dituliskan sebagai daftar, gunakan format berikut.

\begin{enumerate}
  \item Bagaimana merancang solusi untuk permasalahan yang telah diidentifikasi?
  \item Bagaimana mengevaluasi solusi tersebut menggunakan ukuran yang relevan?
\end{enumerate}

\section{Batasan Masalah}

Batasan masalah menjelaskan ruang lingkup data, pengguna, teknologi, lokasi, waktu, atau fitur yang termasuk dan tidak termasuk dalam tugas akhir.

\section{Tujuan Penelitian}

Tujuan harus menjawab rumusan masalah dan menggunakan kata kerja yang dapat dievaluasi, misalnya merancang, mengimplementasikan, membandingkan, atau mengukur.

\section{Manfaat Penelitian}

Jelaskan manfaat akademik dan manfaat praktis yang diharapkan dari hasil tugas akhir.

\section{Sistematika Penulisan}

Jelaskan isi setiap bab secara singkat agar pembaca memahami alur keseluruhan laporan.
